\documentclass[stu,12pt,nolmodern]{apa7}

\usepackage[T1]{fontenc}
\usepackage[utf8]{inputenc}
\usepackage{mathptmx}
\usepackage[spanish,es-nodecimaldot]{babel}
\usepackage{csquotes}
\usepackage{microtype}
\usepackage{booktabs}
\usepackage{graphicx}
\usepackage[
  backend=biber,
  style=apa,
  sorting=nyt,
  sortcites=true
]{biblatex}

\DeclareLanguageMapping{spanish}{spanish-apa}
\addbibresource{referencias.bib}

\title{GymHub: Sistema web para la gestión integral de gimnasios}
\shorttitle{GymHub}
\author{Nombre del estudiante}
\authorsaffiliations{Nombre de la institución}
\course{Código y nombre del curso}
\professor{Nombre del docente}
\duedate{Fecha de entrega}

\abstract{%
  Este documento presenta GymHub, una aplicación web orientada a la gestión de
  miembros, entrenadores, clases, planes de entrenamiento, progreso y
  facturación de un gimnasio. Sustituya este texto con un resumen de entre 150 y
  250 palabras que indique el problema, el objetivo, el método, los resultados
  principales y la conclusión del proyecto.
}

\keywords{gestión de gimnasios, aplicación web, Django, React}

\begin{document}

\maketitle

\section{Introducción}

La administración de un gimnasio requiere coordinar información de miembros,
entrenadores, clases, planes, asistencia, progreso y facturación. GymHub propone
centralizar estos procesos en una aplicación web y reducir la fragmentación de
la información operativa.

Este documento describe el problema atendido, las decisiones técnicas del
proyecto, el método de desarrollo y los resultados de la implementación. Las
citas y la lista de referencias se generan con el estilo APA de séptima edición
\parencite{apa2020}.

\section{Marco teórico}

El backend de GymHub utiliza Django, un marco de trabajo web cuya documentación
define componentes para modelos, solicitudes HTTP, autenticación y seguridad
\parencite{django_documentacion}. La interfaz utiliza React para construir la
experiencia de usuario mediante componentes \parencite{react_documentacion}.

Amplíe esta sección con antecedentes académicos y conceptos directamente
relacionados con el problema. Cada afirmación tomada de otra fuente debe tener
su cita correspondiente.

\section{Metodología}

Describa el enfoque de desarrollo, el levantamiento de requisitos, la
arquitectura, las tecnologías seleccionadas y la estrategia de pruebas. Indique
qué evidencia se utilizó para verificar cada requisito funcional y no
funcional.

\section{Resultados}

Presente los módulos implementados y los resultados verificables. Incluya
tablas, figuras o métricas solo cuando aporten evidencia al análisis. Toda tabla
y figura debe mencionarse en el texto y tener un título descriptivo.

\section{Conclusiones}

Resuma el grado de cumplimiento de los objetivos, las limitaciones encontradas
y las recomendaciones de trabajo futuro. Las conclusiones deben derivarse de
los resultados presentados y no introducir evidencia nueva.


\printbibliography[title={Referencias}]

\end{document}
